% Modified Template by Jonathan Doucette and Kevin Multani, original by: Jonathan Ward

\documentclass[12pt]{article} 
\usepackage[english]{babel}
\usepackage[utf8]{inputenc}
\usepackage{amsmath} % AMS Math Package
\usepackage{amsthm} % Theorem Formatting
\usepackage{amssymb} % Math symbols such as \mathbb
\usepackage{graphicx} % Allows for eps images
\usepackage{multicol} % Allows for multiple columns
\usepackage[dvips,letterpaper,margin=1in,bottom=1in]{geometry}
\usepackage{hyperref}
\usepackage{parskip} % Removes indentation from paragraphs
\usepackage{xcolor,xspace,soul} % Colour, spacing, and highlighting
\usepackage{mathrsfs}
\usepackage{bm} % For bold math symbols
\usepackage{amscd}
\usepackage[all,cmtip]{xy}
%\usepackage{bbm}
\usepackage{titling}
\usepackage{listing} % for code snippets
%\usepackage{minted} % for code snippets
\usepackage{enumerate}
\usepackage{fancyhdr}
\usepackage[]{physics}
\usepackage[makeroom]{cancel}
\usepackage{pdfpages}
\usepackage[]{mcode}
\usepackage[title]{appendix}

% ***********************************************************
% ********************** BEGIN TITLE PAGE *******************
% ***********************************************************
\newcommand*{\titleGM}{\begingroup % Create the command for including the title page in the document
\hbox{ % Horizontal box
\hspace*{0.2\textwidth} % Whitespace to the left of the title page
\rule{1pt}{\textheight} % Vertical line
\hspace*{0.05\textwidth} % Whitespace between the vertical line and title page text
\parbox[b]{0.75\textwidth}{ % Paragraph box which restricts text to less than the width of the page

{\noindent\Huge\bfseries MATH 521}\\[2\baselineskip] % Title
{\large \textit{Assignment 4}}\\[4\baselineskip]
%{\large \textsc{ Jonathan Doucette }} % Author name

\vspace{0.5\textheight} % Whitespace between the title block and the publisher
{\noindent \today }\\[\baselineskip]
%{\noindent Student Number: 35298124 }\\[\baselineskip] 
%{\noindent \footnotesize All problems below are \textit{Copyright \copyright \space 2018 Timm Treskatis. All Rights Reserved.} }\\[\baselineskip]
} % end parbox
} % end hbox
\endgroup}

 % Sets margins and page size
\pagestyle{fancy} 

%\lhead{Jonathan Doucette}
\rhead{\today}
\rfoot{Page \thepage}
\cfoot{}

\makeatletter % Need for anything that contains an @ command 
\renewcommand{\maketitle} % Redefine maketitle to conserve space
{ \begingroup \vskip 10pt \begin{center} \Huge {\bf \@title}
\vskip 10pt \large \@author \hskip 20pt \@date \end{center}
  \vskip 10pt \endgroup \setcounter{footnote}{0} }
\makeatother % End of region containing @ commands

% ***********************************************************
% ********************** END TITLE PAGE *********************
% ***********************************************************

% ***********************************************************
% ********************** BEGIN NEW COMMANDS *****************
% ***********************************************************

\renewcommand{\labelenumi}{(\alph{enumi})} % Use letters for enumerate
\let\vaccent=\v % rename builtin command \v{} to \vaccent{}

%% MISC
\newcommand{\ab}[1]{\left| #1 \right|} % for absolute value
\newcommand{\avg}[1]{\left< #1 \right>} % for average
\let\underdot=\d % rename builtin command \d{} to \underdot{}
\let\baraccent=\= % rename builtin command \= to \baraccent
\renewcommand{\=}[1]{\stackrel{#1}{=}} % for putting numbers above =
\providecommand{\fr}{\frac}
\providecommand{\RR}{\mathbb{R}}
\providecommand{\CC}{\mathbb{C}}
\providecommand{\NN}{\mathbb{N}}
\providecommand{\e}{\epsilon}
\DeclareMathOperator{\di}{d\!}
\newcommand*\ieval[3]{\left.#1\right\rvert_{#2}^{#3}}

%% Vectors
\renewcommand{\v}[1]{\ensuremath{\mathbf{#1}}} 
\newcommand{\gv}[1]{\ensuremath{\mbox{\boldmath$ #1 $}}} % for vectors of Greek letters
\newcommand{\uv}[1]{\ensuremath{\mathbf{\hat{#1}}}} % for unit vector
\providecommand{\wave}[1]{\v{\tilde{#1}}}

%% DERIVATIVES
\renewcommand{\d}[2]{\frac{d #1}{d #2}} % for derivatives
\newcommand{\dubd}[2]{\frac{d^2 #1}{d #2^2}} % for double derivatives
\newcommand{\pd}[2]{\frac{\partial #1}{\partial #2}} % for partial derivatives
\newcommand{\pdd}[2]{\frac{\partial^2 #1}{\partial #2^2}} % for double partial derivatives

%% Operators
\newcommand{\Gradient}{\ensuremath{\mbox{\boldmath$\nabla$}}} % gradient

%% Text
\newcommand{\mathcolorbox}[2]{\colorbox{#1}{$\displaystyle #2$}}
\newcommand{\hlfancy}[3]{\textcolor{#1}{\sethlcolor{#2}\hl{#3}}}
\newcommand{\TODO}[1]{\hlfancy{red}{yellow}{\textbf{TODO: #1}}}

%% Code
%\newcommand{\code}[1]{\mintinline{C}{#1}}
%\newcommand{\code}[1]{\texttt{#1}}
\newcommand{\code}[1]{\lstinline[columns=fixed]{#1}}
\newcommand{\includecode}[1]{\lstinputlisting{#1}}

% ***********************************************************
% ********************** END NEW COMMANDS *******************
% ***********************************************************

% ***********************************************************
% ********************** BEGIN NEW ENVS *********************
% ***********************************************************

% Theorem
\newenvironment{theorem}[2][Theorem]{\begin{trivlist}
\item[\hskip \labelsep {\bfseries #1}\hskip \labelsep {\bfseries #2.}]}{\end{trivlist}}
% Lemma
\newenvironment{lemma}[2][Lemma]{\begin{trivlist}
\item[\hskip \labelsep {\bfseries #1}\hskip \labelsep {\bfseries #2.}]}{\end{trivlist}}
% Corollary
\newenvironment{corollary}[2][Corollary]{\begin{trivlist}
\item[\hskip \labelsep {\bfseries #1}\hskip \labelsep {\bfseries #2.}]}{\end{trivlist}}

% Exercise
\newenvironment{exercise}[2][Exercise]{\begin{trivlist}
\item[\hskip \labelsep {\bfseries #1}\hskip \labelsep {\bfseries #2.}]}{\end{trivlist}}
% Problem
\newenvironment{problem}[2][Problem]{\begin{trivlist}
\item[\hskip \labelsep {\bfseries #1}\hskip \labelsep {\bfseries #2.}]}{\end{trivlist}}
% Question
\newenvironment{question}[2][Question]{\begin{trivlist}
\item[\hskip \labelsep {\bfseries #1}\hskip \labelsep {\bfseries #2.}]}{\end{trivlist}}
% Solution
\newenvironment{solution}{\begin{proof}[Solution]}{\end{proof}}

% Afterword
\newenvironment{afterword}[2][Appendix]{\begin{trivlist}
\item[\hskip \labelsep {\bfseries #1}\hskip \labelsep {\bfseries #2}]}{\end{trivlist}}

% ***********************************************************
% ********************** END NEW ENVS ***********************
% ***********************************************************

% ***********************************************************
% ********************** END TITLEPAGE **********************
% ***********************************************************

%\documentclass[10pt,letterpaper]{scrartcl}
%\usepackage{amsfonts,amsmath,amssymb,braket,xcolor,dsfont,enumerate,fontawesome,graphicx}
\usepackage{amsfonts,amsmath,amssymb,braket,xcolor,enumerate,graphicx}
\usepackage[hidelinks]{hyperref}
\usepackage{listings,multicol,mathtools,textcomp,tikz,pgfplots,wrapfig}
%\usepackage[inner=2cm,outer=2cm,top=2cm,bottom=2cm]{geometry}
\usepackage{tabularx}
\usepackage{booktabs}
\usetikzlibrary{arrows}
%\pgfplotsset{compat=1.12}

\pagestyle{empty}
\setlength{\parindent}{0pt}
\setlength{\parskip}{6pt}

%\newcommand{\dx}{\;\mathrm{d}x}

\begin{document}

\begin{minipage}{.2\textwidth}
\includegraphics[width=42pt]{ubc-logo.png}
\end{minipage}
\hfill
\begin{minipage}{.75\textwidth}
\setlength{\parskip}{6pt}
\begin{flushright}
{\sffamily
\textbf{MATH521}\\
\textbf{Numerical Analysis of Partial Differential Equations}

Winter 2017/18, Term 2\\
%Timm Treskatis
}
\end{flushright}
\end{minipage}

\section*{Homework Assignment 9}

%Please submit the following files as indicated below: \hfill \faFileCodeO \: source code \hfill \faFilePdfO \: PDF file \hfill \faFilePictureO \: image file \hfill \faFileMovieO \: video file

\paragraph*{Question 1 $\vert$ 1 mark}% $\vert$ \faFilePdfO}

We consider the initial boundary value problem for the heat equation
\begin{equation}\tag{H}\label{eq:heat}
\begin{aligned}
\partial_t u(t) - a \Delta u(t) &= f(t) && \text{in } Q = ]0,T[ \times \Omega\\
u(0) &= u_0 && \text{in } \Omega\\
\frac{\partial u}{\partial n} &= 0 && \text{on } \Sigma = ]0,T[ \times \partial \Omega
\end{aligned}
\end{equation}
where $u$ is a temperature field, $u_0$ an initial temperature distribution, the diffusion-like parameter $a > 0$ the heat conductivity of the material, $f$ a source term e.g. due to thermal radiation and $T>0$ a final time. The homogeneous \textsc{Neumann} boundary conditions mean that the domain $\Omega$ is perfectly insulated so that no thermal energy is radiated into the environment.

The $\theta$-method is a class of \textsc{Runge-Kutta} schemes for integrating ODEs of the form
\begin{equation*}
\dot{U} = F(t,U)
\end{equation*}
by using the iteration
\begin{equation*}
U_+ = U_\circ + \Delta t \left( \theta F(t_+,U_+) + (1-\theta) F(t_\circ,U_\circ) \right).
\end{equation*}
The parameter $\theta \in [0,1]$ can be interpreted as the `degree of implicitness', since $\theta = 0$ gives the forward \textsc{Euler} method, $\theta = \frac{1}{2}$ the \textsc{Crank-Nicolson} method (aka implicit trapezium rule in the ODE context) and $\theta = 1$ the backward \textsc{Euler} method.

For the discretisation in space, we apply linear finite elements. Both the time step size $\Delta t$ and the spatial triangulation $\mathcal{T}^h$ are fixed.

Show that in this setting both the method of lines and \textsc{Rothe}'s method lead to the same discrete problems
\begin{equation*}
\left( M^h + \theta \Delta t a K^h\right) \vec{u}^h_+ = \left( M^h - (1-\theta)\Delta t a K^h\right) \vec{u}^h_\circ + \Delta t \left( \theta \vec{f}^h_+ + (1-\theta) \vec{f}^h_\circ \right).
\end{equation*}

You don't have to include any details about the components of the discrete vectors and matrices. We all know what they are!

% ------------------------------------------------------------ %
% Q1: Solution
% ------------------------------------------------------------ %
\vspace{0.5cm}
\begin{solution}

\textbf{Method of lines.} First, we semi-discretise the problem \eqref{eq:heat} in space using linear finite elements to obtain the weak form of the method of lines
\begin{equation}\label{eq:heat-MOL-weak}
\begin{aligned}
M^h \partial_t \vec{u}^h + a K^h \vec{u}^h &= \vec{f}^h && \text{in } Q^h = ]0,T[ \times \Omega^h\\
\vec{u}^h(0) &= \vec{u}^h_0 && \text{in } \Omega^h.
\end{aligned}
\end{equation}
If we now write $U = \vec{u}^h$ and
$F(t,\vec{u}^h) = (M^h)^{-1}( \vec{f}^h - a K^h \vec{u}^h )$, we can write the $\theta$-method scheme for the initial value problem \eqref{eq:heat-MOL-weak} as
\begin{equation*}
\vec{u}^h_+ = \vec{u}^h_\circ + \Delta t
\left( \theta \, (M^h)^{-1}( \vec{f}^h_+ - a K^h \vec{u}^h_+ ) + (1-\theta) (M^h)^{-1}( \vec{f}^h_\circ - a K^h \vec{u}^h_\circ ) \right),
\end{equation*}
or, multiplying through by $M^h$ and rearranging, we have
\begin{equation*}
\left( M^h + \theta \Delta t a K^h \right) \vec{u}^h_+ =
\left( M^h - (1-\theta) \Delta t a K^h \right) \vec{u}^h_\circ +
\Delta t \left( \theta \vec{f}^h_+ + (1-\theta) \vec{f}^h_\circ \right)
\end{equation*}
as desired.

\pagebreak
\textbf{Rothe's method.} First, we semi-discretise equation \eqref{eq:heat} in time using the $\theta$-method by writing $U = u$ and $F(t,u) = f + a \Delta u$.
Since the time step $\Delta t$ is fixed, the resulting scheme is given by
\begin{equation*}
u_+ = u_\circ + \Delta t \Big[ \theta (f_+ + a \Delta u_+) + (1-\theta) (f_\circ + a \Delta u_\circ) \Big].
\end{equation*}
Rearranging, we have
\begin{equation*}
(1 - \theta \Delta t a) \, u_+ =
(1 - (1-\theta) \Delta t a ) \, u_\circ + 
\Delta t \left( \theta f_+ + (1-\theta) f_\circ \right).
\end{equation*}
Now, using the fact that the spatial triangulation $\mathcal{T}^h$ is invariant with time, we can safely multiply both sides through using test functions from the same linear finite element  space and integrate in the standard way to obtain 
\begin{equation*}
\left( M^h + \theta \Delta t a K^h \right) \vec{u}^h_+ =
\left( M^h - (1-\theta) \Delta t a K^h \right) \vec{u}^h_\circ +
\Delta t \left( \theta \vec{f}^h_+ + (1-\theta) \vec{f}^h_\circ \right)
\end{equation*}
as desired.

\end{solution}

\newpage
\paragraph*{Question 2 $\vert$ 4 marks}

\begin{enumerate}[(a)]
\item% \faFileMovieO \:
The \textsf{FEniCS} scridpt \texttt{hw9.py} implements the backward \textsc{Euler} method for Problem \eqref{eq:heat}. Starting from room temperature ($u_0 \equiv 20$), the bottom left corner of a metal piece $\Omega$ with conductivity parameter $a=0.1$ is held over a flame for one second, then the flame is extinguished. This is modelled by
\begin{equation*}
f(t,x) =
\begin{dcases}
200e^{-5x_1^2 - 2x_2^2} & t \leq 1\\
0 & t > 1
\end{dcases}
\end{equation*}
Complete the missing commands to compute the evolution of the temperature field over the first five seconds using a time step size of $\Delta t = 10^{-2}$.

Save your results as a video, using a frame rate such that the video time is equal to the physical time. You don't have to submit any other files for this part of Question 2.

\emph{Hint:} Open the PVD-file in \textsf{ParaView}, click the `Apply'-button, select a reasonable colour map
%\begin{center}
%\includegraphics[scale=1]{colourmap.png}
%\end{center}
and then re-scale the colour values
%\begin{center}
%\includegraphics[scale=1]{rescale.png}
%\end{center}
to the range $[20,160]$. Use the same range for the following questions, too.

% ------------------------------------------------------------ %
% Q2(a): Solution
% ------------------------------------------------------------ %
\vspace{0.5cm}
\begin{solution}

\TODO{see attached video}

\end{solution}

\newpage
\item% \faFileCodeO \:
Generalise this script to implement the $\theta$-method from Question 1. Check whether setting $\theta = 1$ still gives you the same results. Using the same parameters as in Question 2(a), solve the problem with the \textsc{Crank-Nicolson} method and the forward \textsc{Euler} method. What do you observe?

\vspace{6cm}

\item% \faFilePdfO \:
Solve the problem with the forward \textsc{Euler} method again up to time $T=0.1$, once with $\Delta t = 1.25 \times 10^{-4}$ and once with $\Delta t = 10^{-4}$. Explain your observations, using the relevant terminology.

\vspace{6cm}

\item% \faFilePdfO \:
If you could choose among the time-stepping schemes considered above and other time-discretisation schemes you know of, which one would you choose for this problem? First formulate some reasonable objectives for your simulation, then explain which method attains these to the greatest extent.
\end{enumerate}
\newpage

\mbox{}

\vfill

\paragraph*{Your Learning Progress}% $\vert$ \faFilePdfO}

What is the one most important thing that you have learnt from this assignment?

\vspace*{3mm}
\hrulefill

\vspace*{3mm}
\hrulefill

Any new discoveries or achievements towards the objectives of your course project?

\vspace*{3mm}
\hrulefill

\vspace*{3mm}
\hrulefill

What is the most substantial new insight that you have gained from this course this week? Any \emph{aha moment}?

\vspace*{3mm}
\hrulefill

\vspace*{3mm}
\hrulefill

\end{document}
